\part{Tutorial}

Each chapter here is built on a worked example that ships with the code, in
\file{examples/}. They are tiny --- 24\,$\times$\,24\,$\times$\,6 voxels,
seconds on a laptop --- because the point is to show what a switch does, not to
be physically interesting.

Every chapter has the same shape: what the case is, what to run, what you
should see, and what to change to learn something. Four of the cases have an
answer that does not come from this code, and those are flagged: they are the
ones worth trusting.

% =====================================================================
\chapter{Transport with no reaction}
\label{ch:transport}

\emph{Examples 01 and 02.}

\section{The point}

Before adding chemistry, be sure transport is right. These two cases have no
reactions at all, and the second has an answer you can check on paper.

\section{Example 01: flow and a tracer}

A tracer is held at 1 on the inlet face and 0 on the outlet, and reacts with
nothing. \tg{Peclet} is 1, so advection and diffusion are comparable.

\textbf{Expect:} the tracer sweeping left to right; a profile that falls
monotonically with $x$ and never rises; no \code{[NEG!]} warnings; no
dependence on $z$.

\section{Example 02: diffusion only}

The same case with \tg{Peclet} set to 0, which switches the flow solver off
entirely.

\begin{cbnote}
\textbf{\color{cbTeal}This one has a right answer.}\quad
With no flow, no reaction, one face held at 1 and the other at 0, the steady
state is the solution of $\nabla^2 C = 0$ with those boundaries --- a
\textbf{straight line} from 1 to 0 along $x$.

Plot the centreline profile. If it is straight, your diffusion coefficient,
your boundary conditions and your relaxation time are all doing what you think.
If it curves, one of them is not, and it is far easier to find out here than in
a case with biology in it.
\end{cbnote}

\section{What to change}

\begin{center}
\begin{tabular}{@{}p{4.6cm}p{10.2cm}@{}}
\toprule
\textbf{Change} & \textbf{What it tells you}\\
\midrule
\tg{Peclet} 0 $\rightarrow$ 1 $\rightarrow$ 10 & how the balance of advection and diffusion moves the front. At high P\'eclet the front sharpens and the transport solver needs smaller steps.\\
\tg{tau} away from 0.8 & the lattice-Boltzmann relaxation time. Values far from 1 are less stable; below about 0.55 the scheme misbehaves.\\
right boundary Dirichlet $\rightarrow$ Neumann & a closed end: the tracer accumulates instead of leaving.\\
\bottomrule
\end{tabular}
\end{center}

\warn{If example 02 is not straight, stop here.}{Every later chapter builds on
transport being correct. A curved profile with no reaction means something is
firing that should not be, or a boundary is not being held, and both will
corrupt everything downstream while looking plausible.}

% =====================================================================
\chapter{Abiotic kinetics}
\label{ch:abiotic}

\emph{Example 03.}

\section{The point}

Chemistry with no biology. Two reactants enter from opposite faces, meet, and
make a product:
\[
  \mathrm{A} + \mathrm{B} \longrightarrow \mathrm{C},
  \qquad \text{rate} = k\,[\mathrm{A}][\mathrm{B}]
\]

\section{Where the rate law lives}

In \file{defineAbioticKinetics.hh}, in the example's own folder. The whole
function is this:

\begin{lstlisting}[style=cpp]
void defineAbioticRxnKinetics(std::vector<double> C,
                              std::vector<double>& subsR,
                              plb::plint mask)
{
    for (std::size_t i = 0; i < subsR.size(); ++i) subsR[i] = 0.0;
    if (C.size() < 3 || subsR.size() < 3) return;
    if (mask < 2) return;              // nothing happens in solid or wall voxels

    const double A = std::max(C[0], 0.0);
    const double B = std::max(C[1], 0.0);
    const double R = k_AB * A * B;     // mol/L/s

    subsR[0] = -R;
    subsR[1] = -R;
    subsR[2] = +R;
}
\end{lstlisting}

Three things in that snippet are conventions you will reuse everywhere:

\begin{itemize}[leftmargin=1.2em]
\item \textbf{\code{C} is indexed in \file{CompLaB.xml} order.} \code{C[0]} is
\tg{substrate0}. There is no lookup by name.
\item \textbf{\code{subsR} is a rate, per second}, positive for produced and
negative for consumed. The solver multiplies by the time step.
\item \textbf{\code{std::max(x, 0.0)}} floors the tiny negative values the
lattice-Boltzmann solver can produce near steep fronts. Without it a rate law
can amplify numerical noise into a real instability.
\end{itemize}

\section{What you should see}

A peak in C near the middle, with A and B drawn down around it. With
\tg{Peclet} at 0 the front is stationary; turn the flow on and it moves
downstream.

\begin{cbnote}
\textbf{\color{cbTeal}This one has a right answer too.}\quad
One A plus one B makes exactly one C, so summed over the domain
\[
  \Delta[\mathrm{A}] + \Delta[\mathrm{C}] = 0
  \qquad\text{and}\qquad
  \Delta[\mathrm{B}] + \Delta[\mathrm{C}] = 0
\]
at every step. Any drift in those sums is a bug, not chemistry.
\end{cbnote}

% =====================================================================
\chapter{Aqueous speciation}
\label{ch:equilibrium}

\emph{Example 04.}

\section{The point}

Some chemistry is fast enough to treat as instantaneous equilibrium rather than
as a rate. CompLB3D solves a components-and-species tableau in every pore
voxel, at every step.

The example uses the smallest tableau that is still real chemistry: the
carbonate system, four species over two components.

\begin{center}
\begin{tabular}{@{}lccl@{}}
\toprule
\textbf{species} & \textbf{HCO$_3$} & \textbf{H} & \textbf{log\,K}\\
\midrule
HCO$_3$ \emph{(component)}  & 1 & 0 & 0\\
H \emph{(component)}        & 0 & 1 & 0\\
CO$_3$                      & 1 & $-1$ & $-10.33$\\
H$_2$CO$_3$                 & 1 & 1 & $6.35$\\
\bottomrule
\end{tabular}
\end{center}

In XML:

\begin{lstlisting}[style=xml]
<equilibrium>
    <enabled>true</enabled>
    <components>HCO3 H</components>
    <stoichiometry>
        <species0>1  0</species0>
        <species1>0  1</species1>
        <species2>1 -1</species2>
        <species3>1  1</species3>
    </stoichiometry>
    <logK>
        <species0>0.</species0>     <species1>0.</species1>
        <species2>-10.33</species2> <species3>6.35</species3>
    </logK>
</equilibrium>
\end{lstlisting}

\section{The rules of the tableau}

\begin{itemize}[leftmargin=1.2em]
\item \tg{components} names master species, and those names must match
\tg{name\_of\_substrates} exactly.
\item Each \tg{speciesN} row has one number per component, and \code{N} is the
substrate index.
\item A component's own row is the identity row with log\,K of 0.
\item \textbf{A row of zeros means the solver ignores that species.} That is
how minerals and biomass stay out of the tableau while still being substrates.
\end{itemize}

\section{What you should see}

The banner reports the tableau it parsed --- 2 components, 4 species. Acid
enters from the left, so H$_2$CO$_3$ should dominate near the inlet and CO$_3$
away from it.

\begin{cbnote}
\textbf{\color{cbTeal}Check the totals.}\quad
Total carbonate, summed over HCO$_3$, CO$_3$ and H$_2$CO$_3$, must change only
through transport. Speciation moves mass \emph{between} species; it never
creates or destroys a component.
\end{cbnote}

\warn{This is the expensive part.}{Speciation costs more than everything else
in the code combined. Even on 3456 voxels you will see it in the step time.
Leave it off until you need it, and if you do need it, expect the equilibrium
solve to dominate your budget.}

\section{Known limits, stated plainly}

Activities are ideal --- no Debye--H\"uckel or Davies correction. There is no
temperature dependence, no gas phase, no redox couple and no mineral saturation
index. The log\,K values you supply are therefore conditional constants at your
own ionic strength and 25\,\textdegree C. Say so in your methods section.

% =====================================================================
\chapter{Biomass and the three solvers}
\label{ch:biomass}

\emph{Examples 05, 06 and 07 --- the same case, one line different.}

\section{The point}

Biomass has to go somewhere when a colony grows. CompLB3D offers three ways of
moving it, chosen per organism with \tg{solver\_type}. Examples 05 to 07 are
identical apart from that one tag, so running all three shows you exactly what
the choice costs.

\section{The three}

\begin{center}
\begin{tabular}{@{}p{1.5cm}p{6.2cm}p{6.9cm}@{}}
\toprule
& \textbf{How it moves biomass} & \textbf{Use it when}\\
\midrule
\textbf{CA} & cellular automaton: biomass over the maximum density spills into neighbouring voxels & you want a biofilm that spreads by pushing, which is how biofilms actually grow\\
\textbf{FD} & finite difference: biomass diffuses with its own coefficient & you want smooth spreading and a diffusivity you can justify\\
\textbf{LBM} & biomass gets its own advection--diffusion lattice & biomass genuinely moves with the flow\\
\bottomrule
\end{tabular}
\end{center}

\begin{center}
\begin{tikzpicture}[font=\footnotesize,
  n/.style={draw=cbRule,fill=cbSand,rounded corners=2pt,minimum width=2.9cm,
            minimum height=0.8cm,align=center,line width=0.7pt},
  d/.style={draw=cbTeal,fill=white,diamond,aspect=2.6,inner sep=1.5pt,
            align=center,line width=0.7pt},
  a/.style={-{Latex[length=2mm]},cbGrey,line width=0.7pt}]
\node[d] (q1) at (0,0) {does biomass move\\ with the water?};
\node[n] (lbm) at (5.6,0.9) {\textbf{LBM}};
\node[d] (q2) at (0,-2.2) {is it a biofilm that\\ pushes as it grows?};
\node[n] (ca) at (5.6,-1.6) {\textbf{CA}};
\node[n] (fd) at (5.6,-3.0) {\textbf{FD}};
\draw[a] (q1) -- node[above,font=\tiny]{yes} (lbm);
\draw[a] (q1) -- node[left,font=\tiny]{no} (q2);
\draw[a] (q2) -- node[above,font=\tiny]{yes} (ca);
\draw[a] (q2) -- node[below,font=\tiny]{no} (fd);
\end{tikzpicture}
\end{center}

\section{What each one demands of the input file}

\begin{center}
\begin{tabular}{@{}p{1.4cm}p{13.4cm}@{}}
\toprule
\textbf{CA} & needs \tg{viscosity\_ratio\_in\_biofilm}, \tg{thrd\_biofilm\_fraction} and \tg{CA\_method}. \tg{CA\_method} \code{fraction} shares the excess in proportion to the space free in each neighbour.\\
\textbf{FD} & needs \tg{biomass\_diffusion\_coefficients}. \textbf{Rejected for planktonic organisms} --- it requires a \tg{microbeN} entry under \tg{material\_numbers}.\\
\textbf{LBM} & needs \tg{biomass\_diffusion\_coefficients}. Costs one extra lattice per organism, which is the most expensive of the three in memory.\\
\bottomrule
\end{tabular}
\end{center}

\section{The rate law}

\file{defineKinetics.hh} in these folders implements Monod growth:
\[
  \text{growth} = \mu_{\max}\,\frac{S}{K_s+S}\,B,
  \qquad
  \frac{dS}{dt} = -\frac{\text{growth}}{Y},
  \qquad
  \frac{dB}{dt} = \text{growth} - k_d B
\]

\warn{\tg{half\_saturation\_constants} in the XML does not drive this.}{That
tag feeds the flux balance and surrogate paths, which build their own uptake
bounds from it. A hand-written kinetics file takes its parameters from its own
table, in the header, next to the equation that uses them. If you change $K_s$
in the XML and nothing happens, this is why.}

\section{What you should see}

Biomass rising from the seeded patch and then levelling off as growth and decay
balance. The donor drawn down around the colony and replenished from the
boundary. A product plume spreading from the colony. No \code{[NEG!]} warnings.

Run 05, 06 and 07 and compare: the biomass \emph{totals} should be close, the
spatial \emph{patterns} different. That is the honest way to see what the
solver choice actually buys you.

% =====================================================================
\chapter{More than one organism}
\label{ch:twomicrobes}

\emph{Example 08.}

\section{The point}

\tg{solver\_type} and \tg{reaction\_type} are per organism, not global. Example
08 has two microbes on opposite walls sharing one donor, one using the cellular
automaton and the other lattice Boltzmann, in the same run.

They differ in kinetics too: Bug1 has the higher $\mu_{\max}$ and the higher
$K_s$; Bug2 the reverse. So Bug1 wins where the donor is plentiful and Bug2
where it is scarce --- the classic $r$ versus $K$ trade-off. Bug1 sits nearer
the inlet, so it should end up ahead.

\section{What you should see}

Both populations growing, then their \emph{ratio} going flat. A donor shadow
between the two colonies.

\section{What to change}

Swap the two \tg{solver\_type} lines and rerun. The converged ratio should
barely move: the solver decides how biomass spreads, not how much of it there
is. If it moves a lot, that is worth understanding before you trust either
solver on a real problem.

% =====================================================================
\chapter{Flux balance analysis}
\label{ch:fba}

\emph{Examples 09 and 10.}

\section{The point}

Instead of writing a rate law, hand the organism a genome-scale metabolic model
and let a linear program work out how fast it can grow on what is locally
available. CompLB3D solves that program in every voxel, at every step.

\section{The toy model, solved by hand}

Example 09 uses a model small enough to check on paper --- three metabolites,
four reactions:

\begin{center}
\begin{tabular}{@{}clll@{}}
\toprule
\textbf{index} & \textbf{reaction} & \textbf{written as} & \textbf{meaning}\\
\midrule
0 & EX\_S & S$_c$ $\rightarrow$ & negative flux imports the donor\\
1 & EX\_O & O$_c$ $\rightarrow$ & negative flux imports the acceptor\\
2 & EX\_P & P$_c$ $\rightarrow$ & positive flux exports the product\\
3 & BIO   & S$_c$ + 0.5\,O$_c$ $\rightarrow$ P$_c$ & the objective\\
\bottomrule
\end{tabular}
\end{center}

Steady state forces $v_0 = -v_3$ and $v_1 = -\tfrac12 v_3$, so with uptake
bounds $V_S$ and $V_O$:

\begin{cbnote}
\[
  \boxed{\ \text{growth} \;=\; \min\!\left(V_S,\; 2\,V_O\right)\ }
  \qquad\text{where}\qquad
  V = V_{\max}\,\frac{C}{K_c + C}
\]
With \code{fba\_maximum\_uptake\_flux} of 10 and 4, and both concentrations
well above their half-saturation constants, $V_S \to 10$ and $2V_O \to 8$: the
acceptor limits, and growth should settle near \textbf{8\,mmol/gDW/h}.

That number comes from the stoichiometry, not from this code. It is the
strongest check in the suite.
\end{cbnote}

Change \tg{fba\_maximum\_uptake\_flux} to \code{4. 10. 0.} and the donor limits
instead, so growth should settle near 4. One line, a second independent check.

\section{How a voxel gets its answer}

\begin{center}
\begin{tikzpicture}[font=\footnotesize,
  s/.style={draw=cbRule,fill=cbSand,rounded corners=2pt,minimum width=3.0cm,
            minimum height=0.95cm,align=center,line width=0.7pt},
  a/.style={-{Latex[length=2mm]},cbGrey,line width=0.7pt}]
\node[s] (r) at (0,0) {\textbf{1. read}\\ $C$ and $B$ here};
\node[s] (b) at (3.9,0) {\textbf{2. bound}\\ $v_{lb}=-V_{\max}\frac{C}{K_c+C}$};
\node[s] (l) at (7.8,0) {\textbf{3. solve}\\ GLPK or COBRApy};
\node[s] (u) at (3.9,-1.8) {\textbf{4. read back}\\ growth, exchange fluxes};
\node[s] (c) at (7.8,-1.8) {\textbf{5. convert}\\ to $\Delta C$, $\Delta B$};
\draw[a] (r) -- (b);
\draw[a] (b) -- (l);
\draw[a] (l.south) -- ++(0,-0.45) -| (u.north);
\draw[a] (u) -- (c);
\end{tikzpicture}
\end{center}

Steps 1, 2, 4 and 5 are identical for GLPK and COBRApy. Only step 3 differs.

\section{GLPK or COBRApy}

\begin{center}
\begin{tabular}{@{}p{2.2cm}p{6.2cm}p{6.2cm}@{}}
\toprule
& \textbf{GLPK} & \textbf{COBRApy}\\
\midrule
how it works & a persistent \code{glp\_prob}; per voxel it only resets column bounds and re-solves from the previous basis & rebuilds the problem through optlang on every \code{optimize()} call, across a Python boundary\\
speed & the fast path & one to two orders of magnitude slower per voxel\\
needs & \code{-DENABLE\_GLPK=ON} & \code{-DENABLE\_COBRAPY=ON}, cobrapy installed, and \file{src/complab3d\_cobrapy.py} present at run time\\
\bottomrule
\end{tabular}
\end{center}

\warn{COBRApy does not read SBML.}{Both solvers eat the same
\file{*\_rGEM\_c.xml} matrix file produced by \file{extractMM.py}. COBRApy
rebuilds a \code{cobra.Model} in memory from that matrix, with anonymous
numbered reactions --- no gene names, no annotations. Its value here is as an
independent second opinion, not as a richer interface.}

\section{Getting your own model in}

\file{extractMM.py} converts a genome-scale model into the matrix file the
solvers read.

\begin{lstlisting}[style=sh]
python3 extractMM.py iAF987.xml -o input/geobacter.xml
python3 extractMM.py iJO1366.mat -o input/ecoli.xml     # COBRA Toolbox .mat
python3 extractMM.py model.json  -o input/model.xml
\end{lstlisting}

Then in \file{CompLaB.xml}:

\begin{lstlisting}[style=xml]
<model_filename>geobacter</model_filename>          <!-- no .xml -->
<exchange_reaction_indices>6 55 418 -1</exchange_reaction_indices>
<fba_maximum_uptake_flux>15. 25. 8. 0.</fba_maximum_uptake_flux>
<substrate_lower_bounds>-15. -25. -8. 0.</substrate_lower_bounds>
<substrate_upper_bounds>60. 60. 60. 0.</substrate_upper_bounds>
<objective_direction>maximize</objective_direction>
<biomass_molar_mass>24.6</biomass_molar_mass>
\end{lstlisting}

\tg{exchange\_reaction\_indices} maps each substrate, in XML order, onto the
reaction in \emph{that organism's} model that exchanges it. Use \code{-1} for a
substrate the organism does not exchange --- a single \code{-1} is what makes a
mutualism obligate.

\note{\tg{biomass\_molar\_mass}}{converts CompLB3D's molar biomass into the
gDW a metabolic model expects; 24.6\,g/mol is the generic
CH$_{1.8}$O$_{0.5}$N$_{0.2}$ formula weight. If your biomass is already in
gDW/L, set it to 1.}

\section{What you should see}

The model loading with its metabolite and reaction counts. The metabolic
summary naming the solver and the number of microbes using it. Growth
approaching the value the stoichiometry predicts. Product accumulating at the
rate biomass is made.

Run examples 09 and 10 and compare. They receive an identical linear program
and share nothing else, so agreement is strong evidence the coupling is right
and disagreement localises the bug to one of the two.

% =====================================================================
\chapter{Surrogate models}
\label{ch:surrogate}

\emph{Example 11.}

\section{The point}

Flux balance analysis is expensive: one linear program per organism, per active
voxel, per step. On a domain of a few million voxels it dominates everything.

A surrogate replaces the program with a small function fitted to it offline:
\[
  \text{uptake rates} \;\longrightarrow\; \text{growth rate}
\]
evaluated in a few hundred multiplications. Speed-ups of 100 to 1000 times are
typical.

\section{The shipped network}

\begin{center}
\begin{tikzpicture}[font=\scriptsize, x=1.25cm, y=0.42cm,
  nn/.style={circle,draw=cbRule,fill=cbSand,inner sep=1.1pt,line width=0.5pt}]
\foreach \i in {0,1} \node[nn] (i\i) at (0, {2-\i*4}) {};
\foreach \l in {1,2,3,4} \foreach \j in {0,...,4}
  \node[nn] (h\l\j) at (\l, {4-\j*2}) {};
\node[nn] (o) at (5,0) {};
\foreach \i in {0,1} \foreach \j in {0,...,4} \draw[cbRule,line width=0.25pt] (i\i)--(h1\j);
\foreach \l [evaluate=\l as \n using int(\l+1)] in {1,2,3}
  \foreach \j in {0,...,4} \foreach \k in {0,...,4}
    \draw[cbRule,line width=0.25pt] (h\l\j)--(h\n\k);
\foreach \j in {0,...,4} \draw[cbRule,line width=0.25pt] (h4\j)--(o);
\node[cbGrey,anchor=north] at (0,-3.2) {2 inputs};
\node[cbGrey,anchor=north] at (2.5,-3.2) {4 hidden layers, 10 tansig each};
\node[cbGrey,anchor=north] at (5,-3.2) {1 linear};
\end{tikzpicture}
\end{center}

The network shipped in \file{surrogateModel.hh} is \emph{Geobacter
metallireducens} on acetate and Fe(III), exported from MATLAB's Deep Learning
Toolbox.

\begin{cbwarn}
\textbf{\color{cbRed}A network is only valid inside the range it was trained
on.}\quad Outside it, a neural network does not fail --- it returns confident
nonsense. The shipped one was trained on

\begin{center}
\begin{tabular}{@{}lr@{ \ to \ }l@{}}
acetate  & 0.00089 & 9.998 \ mmol/gDW/h\\
Fe(III)  & 0.000035 & 0.4997 \ mmol/gDW/h\\
growth   & 0 & 0.0527 \ 1/h\\
\end{tabular}
\end{center}

Ask it for an acetate uptake of 50 and an Fe(III) uptake of 5, well outside the
box, and it cheerfully reports 0.0486\,1/h.

That is why example 11 sets \tg{fba\_maximum\_uptake\_flux} to \code{8. 0.4} ---
safely inside. Raise the second number above 0.5 and the run spends its time
extrapolating.
\end{cbwarn}

Those bounds are not documented anywhere in the network file. They were
recovered by inverting the \code{mapminmax} scaling, which is what
\file{surrogate\_training/inspectSurrogate.py} does:

\begin{lstlisting}[style=sh]
python3 surrogate_training/inspectSurrogate.py surrogateModel.hh
\end{lstlisting}

Run that on any network before trusting it, including one somebody else fitted.

\section{Fitting your own}

The shipped network encodes one organism on two substrates over one range. For
anything else you need to train your own, and Chapter~\ref{ch:trainsurrogate}
covers that end to end.

% =====================================================================
\chapter{Combining metabolism types}
\label{ch:combined}

\emph{Example 12.}

\section{The point}

A metabolic model describes what the network can do; it does not describe
everything a cell does. Maintenance is the standard example: a constant per
biomass drain that yields no growth. A plain flux balance objective omits it,
so the model over-predicts growth at low substrate.

The combined reaction types run both paths and \textbf{add} the rates:

\begin{center}
\code{glpk\_and\_kinetics} \qquad \code{cobrapy\_and\_kinetics} \qquad
\code{surrogate\_and\_kinetics}
\end{center}

\section{The sign discipline}

\begin{cbwarn}
\textbf{\color{cbRed}Do not write \code{bioR} in a kinetics file that runs
alongside FBA.}\quad Growth belongs to the metabolic path. Writing it in both
places double counts it, and the result looks plausible --- just wrong by a
factor that depends on the substrate field.

Example 12's \file{defineKinetics.hh} writes \code{subsR} only, and says so in
a comment. Copy that discipline.
\end{cbwarn}

\section{Two Vmax tags, not one}

A combined organism needs both:

\begin{center}
\begin{tabular}{@{}p{5.4cm}p{9.4cm}@{}}
\toprule
\tg{maximum\_uptake\_flux} & read by the \textbf{kinetics} path, in mol substrate per mol cells per second\\
\tg{fba\_maximum\_uptake\_flux} & read by the \textbf{metabolic} path, in mmol/gDW/h\\
\bottomrule
\end{tabular}
\end{center}

One number cannot satisfy both meanings. If you supply only the shared tag, the
metabolic layer uses it and prints a warning about the units.

\section{What you should see}

The same growth rate as example 09, because maintenance adds no biomass --- but
the donor drawn down faster, by exactly the maintenance term. Diff the two
donor profiles: the difference \emph{is} the maintenance, and it confirms the
paths are adding rather than one silently winning.

% =====================================================================
\chapter{Precipitation and dissolution}
\label{ch:precip}

\emph{Examples 13, 14 and 15.}

\section{The point}

Minerals form, fill pore space and block the flow; acid arrives and eats them
back out again. The geometry is not a fixed input --- it is a state variable.

\section{The voxel life cycle}

\begin{center}
\begin{tikzpicture}[font=\footnotesize,
  st/.style={draw=cbRule,fill=cbSand,rounded corners=2pt,minimum width=2.8cm,
             minimum height=0.9cm,align=center,line width=0.7pt},
  a/.style={-{Latex[length=2mm]},cbGrey,line width=0.7pt}]
\node[st] (open) at (0,0) {\textbf{open pore}\\ mineral $=0$};
\node[st] (fill) at (4.6,0) {\textbf{filling}\\ $0<$ mineral $<$ full};
\node[st] (seal) at (9.2,0) {\textbf{sealed}\\ mineral $\ge$ full};
\draw[a] (open) -- node[above,font=\tiny]{reaction} (fill);
\draw[a] (fill) -- node[above,font=\tiny]{reaches \texttt{max\_precipRho}} (seal);
\draw[a] (seal.south) to[out=-140,in=-40] node[below,font=\tiny]
  {falls below \texttt{reopen\_fraction} $\times$ full} (fill.south);
\draw[a] (fill.north) to[out=140,in=40] node[above,font=\tiny]{dissolved away} (open.north);
\end{tikzpicture}
\end{center}

\section{Example 13: precipitation}

\[ \mathrm{Fe}^{2+} + \mathrm{HS}^- \longrightarrow \mathrm{FeS(s)} \]

Three things make this work:

\begin{enumerate}[leftmargin=1.4em]
\item The mineral substrate is \textbf{\tg{immobile}}. It never advects or
diffuses, so it accumulates exactly where it forms. A mineral that moved would
never seal anything.
\item \tg{max\_precipRho} is the mol/L of a completely full voxel, derived from
the mineral: $1000 / V_m$ with $V_m$ in cm$^3$/mol. Mackinawite is
87.91\,g/mol over 4.30\,g/cm$^3$, so $V_m = 20.44$ and
\tg{max\_precipRho} $= 48.9$.
\item \tg{surface\_only} \code{1} restricts growth to wall-adjacent voxels ---
heterogeneous nucleation, which is the physically right choice. Set it to 0 and
the mineral grows anywhere, sealing pore throats far too readily.
\end{enumerate}

\textbf{Expect:} FeS building in a band across the middle; the porosity in the
log falling in steps, one per update interval; the flow re-solved after each
conversion; and mass conserved --- every mole of FeS removes one Fe$^{2+}$ and
one HS$^-$.

\section{Example 14: dissolution}

\[ \mathrm{CaCO_3(s)} + \mathrm{H}^+ \longrightarrow \mathrm{Ca}^{2+} + \mathrm{HCO_3}^- \]

Dissolution has its own hook, \code{defineDissolutionRate()}, in the same
header. It is handed:

\begin{center}
\begin{tabular}{@{}p{2.2cm}p{12.6cm}@{}}
\toprule
\code{phaseId} & which declared phase this is, numbered from 1 in \tg{phaseN} order\\
\code{C} & the water \textbf{touching} the mineral: the open neighbours' concentrations, averaged. A solid voxel has no water of its own that moves.\\
\code{mineral} & mol/L of mineral left in this voxel\\
\bottomrule
\end{tabular}
\end{center}

and must return \code{mineralR} \textbf{negative}, plus the released solutes in
\code{subsR}, positive. Do not write the mineral's own \code{subsR} entry --- the
solver does that from \code{mineralR}.

\subsection{Why phases must be declared}

Only what appears in a \tg{phaseN} block can dissolve, and that is deliberate.
If the rule were simply ``a solid voxel with no mineral left becomes pore'',
then an inert quartz grain --- which never held any mineral --- is already below
any threshold, and your whole grain pack would dissolve on the first step.
Declaring phases makes quartz safe by construction.

\begin{lstlisting}[style=xml]
<phase0>
    <name>calcite_grain</name>
    <material_number>0</material_number>   <!-- the solid grains -->
    <substrate>2</substrate>               <!-- holds its inventory -->
    <full_density>27.1</full_density>      <!-- mol/L when full -->
    <initial_fill>27.1</initial_fill>      <!-- starts solid -->
</phase0>
\end{lstlisting}

\tg{initial\_fill} is the point of example 14: the grains \emph{start} solid
and are eaten away. Dissolution is not limited to material that precipitated
first. A precipitate phase instead uses \tg{initial\_fill} 0 and
\tg{is\_precipitate} \code{true}.

\textbf{Expect:} the acid eating into the \emph{upstream} faces first; Ca$^{2+}$
leaving downstream; porosity rising in steps, the mirror of example 13.

\section{Example 15: both at once}

Both blocks enabled, on one mineral, so a voxel can seal and later reopen.

\begin{cbnote}
\textbf{\color{cbTeal}This is what \tg{reopen\_fraction} is for.}\quad
With both directions active a voxel can sit near the threshold. A single
threshold makes it flip open and shut every update interval, and every flip
re-solves the flow. \tg{reopen\_fraction} 0.9 puts a gap between sealing at
full density and reopening below $0.9\times$ full, so it settles instead of
chattering.

Set it to 0.99 and rerun: you should see the porosity oscillate, with a flow
solve wasted on each flip. That is the failure the hysteresis prevents.
\end{cbnote}

The two rate laws are \emph{not} inverses. Precipitation is second order in the
dissolved ions; dissolution is first order in acid and in remaining mineral.
They are separate laws that happen to move the same mineral --- which is how the
code has to work, because the speciation solver computes no solubility
products and so there is no saturation state to drive a single reversible rate.
